% GPQA-Diamond macro analysis section for AlphaDiana paper
% Generated from computed results in analyze_tools/data/
% Scripts: compute_gpqa_action_rates.py, compute_gpqa_goal_loss.py,
%          compute_gpqa_oracle.py, analyze_gpqa_tool_quality.py,
%          build_gpqa_subdomain_map.py
% Requires: \usepackage{tcolorbox} in preamble

% Boxed finding claim, conference benchmark style (steel blue)
\newcounter{finding}
\newcommand{\findingbox}[1]{%
  \refstepcounter{finding}%
  \begin{tcolorbox}[
    colback=blue!6!white,
    colframe=blue!45!black,
    boxrule=0.5pt,
    arc=3pt,
    left=8pt, right=8pt, top=6pt, bottom=6pt,
  ]
  \textbf{Finding \thefinding.}\enspace #1
  \end{tcolorbox}
}


\subsection{GPQA-Diamond}
\label{appendix:macro-gpqa}

We evaluate all four systems on the 198-question GPQA-Diamond benchmark using
Qwen3.5-27B.
Table~\ref{tab:gpqa-overview} reports accuracy, output length, and tool-use rates
for each system.
Two terms used throughout this section: a \emph{harness-specific failure} is a
task DirectLLM answers correctly but the harness answers incorrectly; a
\emph{harness-only success} is a task DirectLLM answers incorrectly but the
harness answers correctly.
The \emph{W/C token ratio} is the mean output length on wrong trajectories
divided by the mean on correct trajectories.
The \emph{tool use rate} is the fraction of trajectories containing at least one
tool invocation.

Three contrasts define the landscape.
OpenClaw's wrong trajectories are nearly twelve times longer than its correct
ones, a W/C ratio absent in any other system.
OpenCode shows a 3.6$\times$ gap in tool use rate between correct and wrong
trajectories, pointing to tool activation as an outcome predictor.
ZeroClaw and OpenClaw have comparable harness-only success rates yet differ
by 13 percentage points in harness-specific failures, indicating that the loss
side of the ledger drives most of the accuracy gap.

A fourth contrast emerges from output entropy.
\emph{Mean entropy} is the average per-token log-probability entropy over a
trajectory, computed from model logprobs; higher entropy indicates more
diffuse, exploratory generation while lower entropy indicates repetitive or
over-confident output.
OpenClaw wrong trajectories collapse to mean entropy 0.14 versus 0.25 on
correct ones: repetitive low-diversity generation accompanying the looping
spiral.
ZeroClaw reverses the direction: wrong trajectories carry higher entropy (0.36)
than correct ones (0.25), reflecting genuine deliberative uncertainty that
nonetheless resolves incorrectly.
OpenCode and DirectLLM show no outcome-conditioned entropy gap, so length
rather than diversity distinguishes their wrong trajectories.

\begin{table}[t]
\centering
\small
\caption{%
  System-level overview on GPQA-Diamond ($n=198$, Qwen3.5-27B).
  Tok\,(C) / Tok\,(W) = mean output tokens on correct / wrong trajectories.
  W/C ratio = Tok\,(W) / Tok\,(C).
  Ent\,(C) / Ent\,(W) = mean per-token log-probability entropy on correct / wrong
  trajectories.
  Tool\,\% = tool use rate, split correct / wrong.
  Loss\,\% = harness-specific failure rate over the 159 DirectLLM-correct tasks.
}
\label{tab:gpqa-overview}
\setlength{\tabcolsep}{3pt}
\begin{tabular}{lrrrrrrr}
\toprule
System & Acc\,\% & Tok (C) & Tok (W) & W/C & Ent (C\,/\,W) & Tool\,\% (C\,/\,W) & Loss\,\% \\
\midrule
DirectLLM &  80.3 &  9{,}049 & 19{,}534 & $2.2{\times}$  & 0.32\,/\,0.30 & 0\,/\,0    & N/A  \\
OpenClaw  &  66.2 &  2{,}372 & 28{,}283 & $11.9{\times}$ & 0.25\,/\,0.14 & 22\,/\,7   & 22.6 \\
OpenCode  &  73.2 &  1{,}066 &  3{,}571 & $3.3{\times}$  & 0.30\,/\,0.30 & 55\,/\,15  & 15.1 \\
ZeroClaw  &  77.8 &  2{,}101 &  2{,}338 & $1.1{\times}$  & 0.25\,/\,0.36 & 0\,/\,0    &  9.4 \\
\midrule
Oracle    &  \textbf{88.4} & \multicolumn{6}{c}{(per-task best available harness)} \\
\bottomrule
\end{tabular}
\end{table}

\paragraph{Action Distribution.}
Table~\ref{tab:gpqa-action-rates} reports the fraction of trajectories in each
(harness, outcome) bucket containing at least one event of each action type.
DirectLLM is a degenerate baseline with no scaffolding, so its rows carry no
diagnostic content beyond outcome counts.
Among the three agent harnesses, action profiles diverge substantially.
OpenCode shows the largest outcome-conditioned gap on tool use (55\% correct
vs.\ 15\% wrong).
OpenClaw shows the opposite on plan (16\% correct vs.\ 31\% wrong): wrong
trajectories over-emit planning steps while under-utilizing recovery.
ZeroClaw is structurally flat across outcomes, consistent with its fixed
plan-then-reason-then-answer template leaving no room for outcome-contingent
divergence.

\begin{table}[t]
\centering
\small
\caption{%
  Outcome-conditioned action rates on GPQA-Diamond ($n=198$).
  Each value is the percentage of trajectories in the (harness, outcome) bucket
  that contain at least one event of that action type.
}
\label{tab:gpqa-action-rates}
\setlength{\tabcolsep}{4pt}
\begin{tabular}{llrcccccc}
\toprule
Harness & Outcome & $n$ & Ans.\,\% & Reason\,\% & Tool\,\% & Verify\,\% & Plan\,\% & Recover\,\% \\
\midrule
DirectLLM & correct & 159 & 100 & 100 &  0 &  0 &  0 &  0 \\
DirectLLM & wrong   &  39 & 100 & 100 &  0 &  0 &  0 &  0 \\
\midrule
OpenClaw  & correct & 131 &  99 &  12 & 22 & 57 & 16 & 24 \\
OpenClaw  & wrong   &  67 &  85 &  19 &  7 & 58 & 31 & 10 \\
\midrule
OpenCode  & correct & 145 & 100 &  93 & 55 & 59 & 30 &  0 \\
OpenCode  & wrong   &  53 & 100 &  92 & 15 & 49 & 32 &  0 \\
\midrule
ZeroClaw  & correct & 154 & 100 & 100 &  0 & 60 & 82 &  0 \\
ZeroClaw  & wrong   &  37 & 100 & 100 &  0 & 59 & 81 &  0 \\
\bottomrule
\end{tabular}
\end{table}

\findingbox{Agent loops without a commit deadline produce systematic trajectory
inflation on wrong instances; context exhaustion rather than reasoning failure
becomes the proximate cause of incorrect answers.}

A harness loop without a commit deadline converts deliberation into a liability
on hard instances.
When reasoning converges, OpenClaw's open loop is harmless; when it does not,
the model re-enters planning or verification indefinitely until the context
budget is exhausted.
Table~\ref{tab:gpqa-goal-loss} captures this asymmetry: correct OpenClaw
trajectories are brief, while wrong ones spiral to an average of 28k tokens and
34\% hit the context limit before the model can emit an answer.
ZeroClaw's structurally fixed template prevents re-entry entirely, producing
nearly identical token budgets on correct and wrong trajectories.
DirectLLM's cap events reflect a separate mechanism: a confident but incorrect
chain generated at low entropy, running out of space before a commit decision
rather than cycling.
The same surface symptom (long wrong output) has two distinct causes requiring
different interventions.

\textbf{Case: gpqa\_140 (Organic Chemistry).}
Correct chemical analysis appears early in the OpenClaw trajectory, but the
loop re-enters before any commit, exhausts the context limit mid-sentence, and
produces no extractable answer.

\begin{table}[t]
\centering
\small
\caption{%
  Trajectory-level signals by system and outcome on GPQA-Diamond.
  Mean tokens = average output length.
  Cap\,\% = percentage of trajectories reaching the 32k-token context limit.
  DirectLLM is single-pass; cap events reflect confident overrun, not looping.
}
\label{tab:gpqa-goal-loss}
\setlength{\tabcolsep}{4pt}
\begin{tabular}{llrrr}
\toprule
System & Outcome & $n$ & Mean tokens & Cap\,\% \\
\midrule
DirectLLM & correct & 159 &  9{,}049 &  0 \\
DirectLLM & wrong   &  39 & 19{,}534 & 33 \\
\midrule
OpenClaw  & correct & 131 &  2{,}372 &  0 \\
OpenClaw  & wrong   &  67 & 28{,}283 & 34 \\
\midrule
OpenCode  & correct & 145 &  1{,}066 &  0 \\
OpenCode  & wrong   &  53 &  3{,}571 &  0 \\
\midrule
ZeroClaw  & correct & 154 &  2{,}101 &  0 \\
ZeroClaw  & wrong   &  37 &  2{,}338 &  0 \\
\bottomrule
\end{tabular}
\end{table}

\findingbox{Tool activation is outcome-predictive only when the task affords an
executable path; the same harness and model produce opposite outcome-conditioned
tool-use gaps across structurally different subdomains.}

Tool activation is the strongest outcome-conditioned action signal on this
benchmark, but only when paired with a useful tool result.
Table~\ref{tab:gpqa-tool-quality} shows that restricting to substantive results
(non-empty, no error string, over 100 characters) preserves nearly the full
correct-vs-wrong gap, ruling out trajectories that invoke a tool but discard its
output.
The gap is subdomain-driven rather than model-driven: tasks with tractable
numerical or symbolic structure afford an executable path, where code execution
anchors the answer; tasks with irreducibly verbal reasoning (stereochemistry,
reaction mechanisms) offer no such path, and the deliberation loop from
Finding~1 dominates.

\textbf{Case contrast: gpqa\_167 vs.\ gpqa\_143.}
On a Quantum Mechanics eigenvector problem, OpenCode executes three code calls
and commits correctly in under 1,400 tokens.
On a stereochemistry problem, zero tool calls are made and the model commits to
the wrong option in comparable token length.
Same harness, same model; task structure determines whether the tool path opens.

\begin{table}[t]
\centering
\small
\caption{%
  Tool-use quality by outcome on GPQA-Diamond.
  Any call\,\% = trajectories with $\geq$1 tool invocation.
  Substantive\,\% = trajectories whose tool result is non-empty, longer than
  100 characters, and contains no error string.
  Err/empty\,\% = trajectories with an empty or error-containing result.
  Mean calls = average invocations per trajectory.
  OpenClaw values from normalized trace; Err/empty not tracked for OpenCode (N/A).
}
\label{tab:gpqa-tool-quality}
\setlength{\tabcolsep}{4pt}
\begin{tabular}{llrrrrr}
\toprule
Harness & Outcome & $n$ & Any call\,\% & Subst.\,\% & Err/empty\,\% & Mean calls \\
\midrule
OpenClaw & correct & 131 & 56 & 51 & 23 & 1.34 \\
OpenClaw & wrong   &  67 & 19 & 15 & 12 & 0.69 \\
\midrule
OpenCode & correct & 145 & 55 & 55 & N/A & 0.95 \\
OpenCode & wrong   &  53 & 15 & 15 & N/A & 0.32 \\
\bottomrule
\end{tabular}
\end{table}

\findingbox{The same ``wrong'' label aggregates mechanistically distinct failure
types: extraction failure from context truncation, premature commitment under
budget pressure, and genuine reasoning error each leave different signatures in
output format and trajectory length.}

The ``harness loses'' row of Table~\ref{tab:gpqa-failure-sig} isolates
harness-induced damage from pre-existing model limitations.
\emph{Malformed} means the predicted answer falls outside $\{\mathrm{A,B,C,D}\}$;
\emph{No boxed} means no \texttt{\textbackslash boxed\{$\cdot$\}} token appears
and the extractor cannot recover any answer; \emph{TokRatio} is the harness
token count divided by the DirectLLM token count on the same task.

Each harness exhibits a structurally distinct failure signature in that row.
OpenClaw regressions are predominantly malformed and missing a boxed answer at
a TokRatio above 2: the loop spirals past context and the extractor recovers
nothing.
OpenCode regressions are partially malformed but never missing a boxed answer,
and TokRatio is well below 1: the format is intact, but the answer was reached
at a fraction of DirectLLM's deliberation budget.
ZeroClaw regressions are largely well-formed, short, and structurally normal;
the model committed to the wrong answer without any harness artifact.

\textbf{Case: gpqa\_143 (Organic Chemistry, Cope rearrangement).}
All three failure modes appear on the same task: OpenClaw spirals past context
with no extractable output, OpenCode emits a well-formed but wrong answer with
minimal deliberation, and ZeroClaw outputs a malformed option in 590 tokens,
all labeled ``wrong'' without distinction.

\begin{table}[t]
\centering
\small
\caption{%
  Paired-outcome failure signatures on GPQA-Diamond ($n=198$).
  ``Harness loses'' = DirectLLM correct, harness wrong.
  ``Harness gains'' = DirectLLM wrong, harness correct.
  Malformed\,\% = predicted answer not in $\{\mathrm{A,B,C,D}\}$.
  No boxed\,\% = no \texttt{\textbackslash boxed\{$\cdot$\}} in raw output.
  TokRatio = harness token count / DirectLLM token count on the same task.
}
\label{tab:gpqa-failure-sig}
\setlength{\tabcolsep}{3.5pt}
\begin{tabular}{llrrrr}
\toprule
Harness & Outcome & $n$ & Malformed\,\% & No boxed\,\% & TokRatio \\
\midrule
\multirow{4}{*}{OpenClaw}
 & both correct  & 123 &  0.8 &  0.8 & 0.27 \\
 & harness gains &   8 &  0.0 &  0.0 & 0.30 \\
 & harness loses &  36 & 89.0 & 83.0 & 2.47 \\
 & both wrong    &  31 & 58.0 & 58.0 & 1.02 \\
\midrule
\multirow{4}{*}{OpenCode}
 & both correct  & 135 &  0.0 &  0.0 & 0.13 \\
 & harness gains &  10 &  0.0 &  0.0 & 0.08 \\
 & harness loses &  24 & 63.0 &  0.0 & 0.29 \\
 & both wrong    &  29 & 48.0 &  0.0 & 0.17 \\
\midrule
\multirow{4}{*}{ZeroClaw}
 & both correct  & 141 &  0.0 &  0.0 & 0.25 \\
 & harness gains &  13 &  0.0 &  0.0 & 0.11 \\
 & harness loses &  15 & 20.0 &  0.0 & 0.17 \\
 & both wrong    &  22 & 18.0 &  0.0 & 0.17 \\
\bottomrule
\end{tabular}
\end{table}

\findingbox{Harnesses consistently damage more answers than they recover; because
correct instances outnumber recoverable failures, reducing harness-induced loss
is a more tractable accuracy lever than maximizing recovery.}

The \emph{recovery rate} measures how often a harness answers correctly on tasks
DirectLLM fails; the \emph{loss rate} measures how often it answers incorrectly
on tasks DirectLLM already solves.
Table~\ref{tab:gpqa-oracle} shows that recovery rates across all three harnesses
collectively cover only 41\% of DirectLLM's failures, and individually span a
narrow range.
Loss rates, by contrast, vary by more than a factor of two.
Since DirectLLM-correct tasks outnumber DirectLLM-wrong tasks four to one on
this benchmark, reducing the loss rate produces more accuracy gain than an
equivalent improvement in recovery.
ZeroClaw leads not by recovering more DirectLLM failures but by losing fewer
tasks that DirectLLM already handles, confirming that the loss rate is the
tractable lever for improvement.

Subdomain analysis (Table~\ref{tab:gpqa-subdomain}) shows that losses concentrate
where tool use is unavailable and deliberation loops are most active (Organic
Chemistry), while the only case where a harness materially exceeds DirectLLM
(ZeroClaw on Quantum Mechanics) is a subdomain whose sequential algebraic
structure matches ZeroClaw's fixed template precisely.
Harnesses gain where tasks are decomposable and executable; they lose where
reasoning is irreducibly verbal and multi-hop.

\begin{table}[t]
\centering
\small
\caption{%
  Oracle harness router ceiling on GPQA-Diamond ($n=198$).
  Recovery\,\% = harness-only success rate over the 39 DirectLLM-wrong tasks.
  Loss\,\% = harness-specific failure rate over the 159 DirectLLM-correct tasks.
  Oracle selects, per task, any harness that answers correctly when one exists.
}
\label{tab:gpqa-oracle}
\setlength{\tabcolsep}{4pt}
\begin{tabular}{lrrr}
\toprule
System & Accuracy\,\% & Recovery\,\% & Loss\,\% \\
\midrule
DirectLLM & 80.3 & N/A  & N/A \\
OpenClaw  & 66.2 & 20.5 & 22.6 \\
OpenCode  & 73.2 & 25.6 & 15.1 \\
ZeroClaw  & 77.8 & 33.3 &  9.4 \\
\midrule
Oracle    & \textbf{88.4} & 41.0 & N/A \\
\bottomrule
\end{tabular}
\end{table}

\begin{table}[t]
\centering
\small
\caption{%
  Pass rates by GPQA-Diamond subdomain, sorted by DirectLLM accuracy descending.
}
\label{tab:gpqa-subdomain}
\setlength{\tabcolsep}{4pt}
\begin{tabular}{lrcccc}
\toprule
Subdomain & $n$ & DirectLLM & OpenClaw & OpenCode & ZeroClaw \\
\midrule
High-energy particle physics  & 14 & 100 & 100 & 100 & 100 \\
Electromagnetism \& Photonics &  6 & 100 &  83 &  83 & 100 \\
Condensed Matter Physics      &  1 & 100 & 100 & 100 & 100 \\
Optics and Acoustics          &  1 & 100 & 100 & 100 & 100 \\
Inorganic Chemistry           &  1 & 100 & 100 &   0 & 100 \\
Quantum Mechanics             & 25 & 100 &  68 &  84 & 100 \\
Physics (general)             & 19 &  95 &  79 &  89 &  95 \\
Relativistic Mechanics        &  7 &  86 &  71 &  86 &  86 \\
Chemistry (general)           & 20 &  85 &  75 &  75 &  80 \\
Astrophysics                  & 13 &  85 &  77 &  92 & 100 \\
Molecular Biology             & 15 &  73 &  67 &  73 &  60 \\
Organic Chemistry             & 72 &  64 &  49 &  57 &  66 \\
Genetics                      &  4 &  50 &  50 &  25 &  25 \\
\bottomrule
\end{tabular}
\end{table}

\findingbox{Output entropy is a harness-conditional failure signal: low-entropy-long
trajectories indicate pathological repetition under open loops, while
high-entropy trajectories indicate unresolved deliberation under structured
templates; the two failure regimes require different diagnostic criteria.}

Output entropy, computed from per-token log-probabilities, provides a
trajectory-level signal whose direction is harness-specific.
Table~\ref{tab:gpqa-entropy} shows that OpenClaw wrong trajectories collapse to
half the entropy of correct ones while simultaneously growing to 28k tokens.
The model is not expressing uncertainty; it cycles through low-diversity tokens
as the loop runs.
Partitioning by the 25th-percentile entropy and 75th-percentile token threshold
isolates a low-entropy-long quadrant of 41 tasks at 7\% accuracy: a near-certain
failure predictor for OpenClaw specifically.

ZeroClaw inverts this pattern entirely.
Wrong trajectories carry higher entropy than correct ones (0.36 vs.\ 0.25): the
model explores genuinely under uncertainty and still arrives at the wrong answer.
The low-entropy-long quadrant for ZeroClaw is high-accuracy (93\%); short,
confident trajectories are correct under its structured template, not indicative
of pathological looping.
OpenCode shows a flat entropy profile across both outcomes; length rather than
output diversity separates its correct from wrong trajectories.
Entropy is therefore not a universal failure signal but a harness-conditional one.

\begin{table}[t]
\centering
\small
\caption{%
  Entropy and token statistics by system and outcome on GPQA-Diamond (tasks with
  logprob data; OpenClaw $n=186$, OpenCode $n=193$, ZeroClaw $n=191$,
  DirectLLM $n=198$).
  Mean entropy = mean per-token log-probability entropy over the trajectory.
  Cap\,\% = fraction reaching the 32k-token context limit (all cap events are
  wrong).
}
\label{tab:gpqa-entropy}
\setlength{\tabcolsep}{4pt}
\begin{tabular}{llrrrr}
\toprule
System & Outcome & $n$ & Mean entropy & Mean tokens & Cap\,\% \\
\midrule
DirectLLM & correct & 159 & 0.322 &  9{,}050 &  0 \\
DirectLLM & wrong   &  39 & 0.295 & 19{,}534 & 33 \\
\midrule
OpenClaw  & correct & 128 & 0.254 &  2{,}373 &  0 \\
OpenClaw  & wrong   &  58 & 0.139 & 28{,}284 & 40 \\
\midrule
OpenCode  & correct & 145 & 0.301 &  1{,}067 &  0 \\
OpenCode  & wrong   &  48 & 0.300 &  3{,}244 &  0 \\
\midrule
ZeroClaw  & correct & 154 & 0.248 &  2{,}102 &  0 \\
ZeroClaw  & wrong   &  37 & 0.363 &  2{,}338 &  0 \\
\bottomrule
\end{tabular}
\end{table}

% ===========================================================================
\paragraph{Cognitive-Mode Analysis: the ARM Taxonomy.}
% ===========================================================================

The action-rate analysis tells us \emph{what} operations each harness performs.
The ARM (Agentic Reasoning Mode) taxonomy asks a different question: given a
reasoning step, \emph{what cognitive strategy} does it represent?
We classify every assistant-generated segment of every GPQA-Diamond trajectory
into one of six reasoning modes plus an undifferentiated \texttt{reason}
catch-all:

\begin{itemize}[leftmargin=*, itemsep=1pt]
\item \texttt{PD} (Principle-Based Deduction): derivation from physical laws,
  chemical mechanisms, or mathematical theorems
\item \texttt{SE} (Systematic Elimination): explicit option-by-option comparison
\item \texttt{IC} (Intuitive Commitment): direct assertion without extended
  derivation
\item \texttt{UN} (Uncertainty Navigation): explicit hedging or probabilistic
  language
\item \texttt{RR} (Recovery-Replanning): correction of prior reasoning or
  direction change
\item \texttt{reason}: generic deliberation not triggering any classifier
\end{itemize}

Classification is per-segment (each assistant message = one segment) via
keyword heuristics with domain-specific chemistry, physics, and biology signals.
Arm entropy measures cognitive mode diversity within a trajectory.

% ---------------------------------------------------------------------------
\paragraph{PD over-expression is the strongest single-mode predictor of failure.}
% ---------------------------------------------------------------------------

Table~\ref{tab:gpqa-arm-modes} reports the mean ARM mode rate per
(harness, outcome) cell.

\begin{table}[t]
\centering
\small
\caption{%
  ARM cognitive-mode segment allocation by harness and outcome on GPQA-Diamond.
  Each cell is the mean fraction of segments classified to that mode.
  Arm entropy = mode diversity (higher = more modes used within a trajectory).
  $\Delta$PD = PD(wrong) $-$ PD(correct) in percentage points.
}
\label{tab:gpqa-arm-modes}
\setlength{\tabcolsep}{3pt}
\begin{tabular}{llrrrrrrr}
\toprule
Harness & Out. & PD & reason & IC & UN & RR & SE & Arm ent. \\
\midrule
\multirow{2}{*}{OpenClaw}
  & C & 0.519 & 0.386 & 0.052 & 0.017 & 0.018 & 0.007 & 0.548 \\
  & W & \textbf{0.608} & \textbf{0.291} & \textbf{0.026} & \textbf{0.043} & 0.010 & 0.022 & \textbf{0.395} \\
  & $\Delta$ & +8.9 & $-$9.5 & $-$2.6 & +2.6 & $-$0.8 & +1.5 & $-$0.15 \\
\midrule
\multirow{2}{*}{OpenCode}
  & C & 0.353 & 0.587 & 0.026 & 0.009 & 0.019 & 0.006 & 0.593 \\
  & W & \textbf{0.498} & \textbf{0.450} & 0.000 & \textbf{0.041} & 0.010 & 0.000 & \textbf{0.468} \\
  & $\Delta$ & +14.5 & $-$13.7 & $-$2.6 & +3.2 & $-$0.9 & $-$0.6 & $-$0.13 \\
\midrule
\multirow{2}{*}{ZeroClaw}
  & C & 0.280 & 0.674 & 0.023 & 0.011 & 0.021 & 0.001 & 0.636 \\
  & W & \textbf{0.406} & \textbf{0.504} & 0.011 & \textbf{0.032} & \textbf{0.044} & 0.003 & \textbf{0.812} \\
  & $\Delta$ & +12.6 & $-$17.0 & $-$1.2 & +2.2 & +2.4 & +0.2 & \textbf{+0.18} \\
\bottomrule
\end{tabular}
\end{table}

\findingbox{Principle-based deduction is the dominant cognitive mode on
expert-level science questions, and its over-representation is the strongest
single-mode predictor of failure across all three harness architectures.}

GPQA-Diamond is a hard-science benchmark; all three harnesses spend the
plurality of their reasoning budget in PD.
But the allocation is not outcome-neutral: wrong trajectories allocate 8.9 to
14.5~pp \emph{more} of their segments to PD than correct trajectories do, at
the direct expense of undifferentiated \texttt{reason} and IC.
The PD/IC ratio, measuring the tendency to keep deriving rather than committing,
rises 12\% to 55\% in wrong trajectories across all three harnesses.
The interpretation is not that deduction is harmful, but that failure
trajectories are trapped in a deduction spiral: the model re-derives known
principles without transitioning to elimination or commitment.

Arm entropy direction is harness-specific.
OpenClaw and OpenCode wrong trajectories have \emph{lower} mode diversity than
correct ones ($-$0.15 and $-$0.13 respectively): the model collapses into fewer
modes.
ZeroClaw wrong trajectories have \emph{higher} mode diversity (+0.18): the model
oscillates between PD, reason, UN, and RR without converging.
The same ARM marker (high arm entropy) therefore signals productive exploration
under ZeroClaw but unproductive mode oscillation under OpenClaw.

% ---------------------------------------------------------------------------
\paragraph{PD self-transition creates a near-absorbing cognitive state.}
% ---------------------------------------------------------------------------

Table~\ref{tab:gpqa-pd-transitions} reports ARM transition probabilities
for the PD-related cognitive moves.

\begin{table}[t]
\centering
\small
\caption{%
  PD-related ARM transition probabilities by harness and outcome on
  GPQA-Diamond.
  PD$\to$PD = probability a PD segment is followed by another PD segment.
  reason$\to$PD = probability generic reasoning enters deduction.
  reason$\to$IC = probability reasoning exits to intuitive commitment.
}
\label{tab:gpqa-pd-transitions}
\setlength{\tabcolsep}{4pt}
\begin{tabular}{llrrrr}
\toprule
Harness & Outcome & PD$\to$PD & reason$\to$PD & reason$\to$IC & PD$\to$reason \\
\midrule
\multirow{2}{*}{OpenClaw}
  & C & 0.729 & 0.513 & \textbf{0.356} & 0.248 \\
  & W & \textbf{0.863} & \textbf{0.765} & \textbf{0.091} & 0.345 \\
\cmidrule{1-5}
\multirow{2}{*}{OpenCode}
  & C & 0.251 & 0.301 & \textbf{0.198} & 0.222 \\
  & W & \textbf{0.536} & \textbf{0.525} & \textbf{0.016} & 0.270 \\
\cmidrule{1-5}
\multirow{2}{*}{ZeroClaw}
  & C & 0.215 & 0.165 & 0.079 & 0.163 \\
  & W & 0.241 & 0.182 & 0.083 & 0.182 \\
\bottomrule
\end{tabular}
\end{table}

\findingbox{The PD$\to$PD self-transition probability reaches 0.86 on OpenClaw
wrong trajectories, creating a near-absorbing cognitive state; simultaneously,
the commitment exit path (reason$\to$IC) collapses 4$\times$ to 10$\times$
across open-loop harnesses.}

Three transition signatures distinguish spiraling from non-spiraling harnesses.
First, OpenClaw wrong trajectories show PD$\to$PD at 0.86: once the model enters
deduction under uncertainty, it has an 86\% probability of remaining in
deduction on the next cognitive segment.
Second, the commitment exit path collapses: reason$\to$IC drops from 0.36
(correct) to 0.09 (wrong) on OpenClaw, and from 0.20 to 0.02 on OpenCode.
Third, the re-entry probability from generic reasoning into PD (reason$\to$PD)
rises sharply on wrong trajectories for OpenClaw (0.51 to 0.77) and OpenCode
(0.30 to 0.53), but remains flat for ZeroClaw (0.17 to 0.18): the fixed
template prevents re-entry even when the model is uncertain.

These three probabilities encode the mechanism: a near-absorbing PD state, a
collapsing exit path, and an architecture-dependent re-entry trigger.
A harness that breaks any one of the three prevents the entropy collapse
documented in Finding~\ref{finding} (the most recent Finding~5 on entropy).

% ---------------------------------------------------------------------------
\paragraph{Recovery is efficacy-heterogeneous across paired outcomes.}
% ---------------------------------------------------------------------------

Table~\ref{tab:gpqa-arm-paired} reports ARM mode rates disaggregated by the
four-way paired outcome decomposition from Table~\ref{tab:gpqa-failure-sig}.

\begin{table}[t]
\centering
\small
\caption{%
  ARM mode rates by paired outcome on GPQA-Diamond.
  Arm entropy = cognitive mode diversity.
  Regression = DirectLLM correct, harness wrong.
  Rescue = DirectLLM wrong, harness correct.
}
\label{tab:gpqa-arm-paired}
\setlength{\tabcolsep}{3pt}
\begin{tabular}{llrrrrr}
\toprule
Harness & Paired outcome & RR & Arm ent. & PD & IC & Tokens \\
\midrule
\multirow{4}{*}{OpenClaw}
  & both correct & 0.018 & 0.476 & 0.519 & 0.058 &  1{,}932 \\
  & both wrong   & 0.016 & 0.324 & 0.568 & 0.024 & 22{,}148 \\
  & regression   & 0.005 & 0.261 & 0.643 & 0.028 & 30{,}451 \\
  & rescue       & 0.000 & 0.161 & 0.831 & 0.004 &  3{,}598 \\
\midrule
\multirow{4}{*}{OpenCode}
  & both correct & 0.020 & 0.563 & 0.335 & 0.027 &    967 \\
  & both wrong   & 0.002 & 0.486 & 0.471 & 0.001 &  3{,}251 \\
  & regression   & 0.021 & 0.407 & 0.530 & 0.000 &  3{,}672 \\
  & rescue       & 0.000 & 0.386 & 0.600 & 0.014 &    701 \\
\midrule
\multirow{4}{*}{ZeroClaw}
  & both correct & 0.020 & 0.650 & 0.271 & 0.023 &  1{,}982 \\
  & both wrong   & 0.044 & 0.762 & 0.392 & 0.019 &  2{,}482 \\
  & regression   & 0.044 & 0.825 & 0.427 & 0.000 &  1{,}913 \\
  & rescue       & 0.032 & 0.679 & 0.380 & 0.014 &  2{,}333 \\
\bottomrule
\end{tabular}
\end{table}

\findingbox{Recovery (RR) is not a portable quality signal: it marks productive
self-correction under OpenClaw and OpenCode, but mode oscillation under
ZeroClaw, where wrong trajectories carry the highest RR rates and highest arm
entropy simultaneously.}

Under OpenClaw and OpenCode, RR appears in 9\% to 12\% of correct trajectories
versus 3\% to 4\% of wrong ones: recovery is a success signal.
Under ZeroClaw, the direction inverts: 49\% of wrong trajectories contain RR
versus 24\% of correct ones.
But ZeroClaw wrong trajectories also carry the highest arm entropy of any GPQA
cell (0.81 to 0.83 for both-wrong and regression), while their RR rate is
0.044: the RR signal is embedded in mode oscillation rather than focused
correction.

The quadruple (RR rate, arm entropy, PD rate, paired outcome) jointly
determines whether recovery is productive.
Productive recovery (OpenClaw/OpenCode correct): moderate arm entropy
(0.48 to 0.56), moderate PD (0.34 to 0.52), RR present as a focused correction.
Mode oscillation (ZeroClaw wrong): high arm entropy (0.76 to 0.83), elevated
PD (0.39 to 0.43), RR embedded in mode switching without convergence.
Cognitive collapse (OpenClaw regression): low arm entropy (0.26), very high PD
(0.64), RR near zero: the model cannot recover because it cannot exit deduction.

% ---------------------------------------------------------------------------
\paragraph{Tool-type composition and the intent-execution gap.}
% ---------------------------------------------------------------------------

The canonical action analysis (Table~\ref{tab:gpqa-action-rates}) reports
\emph{whether} a harness calls tools.
Table~\ref{tab:gpqa-tool-composition} decomposes \emph{what kind} of tools and
\emph{whether the calls are substantive}.
OpenClaw provides a general-purpose tool shed (shell, read, search, write,
memory\_read); OpenCode provides a single code-execution environment;
ZeroClaw provides none.

\begin{table}[t]
\centering
\small
\caption{%
  Tool-type composition and call quality by harness and outcome on GPQA-Diamond.
  Tool types from event-log subtype labels.
  Any call = mean tool invocations per trajectory.
  Substantive = mean invocations producing executable code or non-empty output.
  Intent = keyword markers expressing desire to use a tool without necessarily
  executing one.
}
\label{tab:gpqa-tool-composition}
\setlength{\tabcolsep}{3pt}
\begin{tabular}{llrrrrrrr}
\toprule
Harness & Out. & Shell & Read & Search & Other & Any & Subst. & Intent \\
\midrule
\multirow{2}{*}{OpenClaw}
  & C & 67\% & 12\% & 9\% & 12\% & 1.34 & 0.89 &   2.1 \\
  & W & 50\% & 17\% & 17\% & 17\% & 0.69 & 0.28 & \textbf{48.0} \\
\midrule
\multirow{2}{*}{OpenCode}
  & C & \multicolumn{4}{c}{100\% code execution} & 0.95 & 0.95 &   7.8 \\
  & W & \multicolumn{4}{c}{100\% code execution} & 0.32 & 0.32 &  11.4 \\
\midrule
\multirow{2}{*}{ZeroClaw}
  & C & \multicolumn{4}{c}{no tools available} & 0.00 & 0.00 &   6.8 \\
  & W & \multicolumn{4}{c}{no tools available} & 0.00 & 0.00 &   7.2 \\
\bottomrule
\end{tabular}
\end{table}

\findingbox{Tool-intent inflation without tool execution is the lexical
signature of the OpenClaw spiral: intent markers explode 23-fold on wrong
trajectories (2.1 to 48.0) while actual tool calls drop to near zero (0.40 to
0.09).
OpenCode inverts the pattern: every call is substantive, but the call rate
collapses 3.0$\times$ when the model is on an unproductive path.}

Three structural facts emerge.
First, OpenClaw provides multiple tool types (shell, read, search) but the model
cannot convert intent into execution on wrong trajectories.
Correct trajectories average 1.34 tool calls with 67\% substantive execution;
wrong trajectories drop to 0.69 calls with 41\% substantive.
Meanwhile, tool-intent markers explode from 2.1 to 48.0: the model repeatedly
expresses the desire to compute something without executing.

Second, OpenCode's single-tool surface produces an all-or-nothing pattern.
Every tool call is substantive (code execution always produces output), but the
call rate drops from 0.95 (correct) to 0.32 (wrong).
Tool activation is outcome-predictive not because tools cause correctness, but
because the model persists in calling tools only when on a productive path.

Third, ZeroClaw has no tool surface, yet its tool-intent markers (6.8 to 7.2)
are comparable to OpenCode correct trajectories (7.8).
The model expresses approximately the same baseline level of computational
intent regardless of harness; OpenClaw's open loop amplifies this into a
repeating marker without execution, while ZeroClaw's template suppresses it.

% ---------------------------------------------------------------------------
\paragraph{Action motifs distinguish productive from pathological trajectories.}
% ---------------------------------------------------------------------------

Canonical actions $\{\texttt{answer}, \texttt{reason}, \texttt{tool\_use},
\texttt{verify}, \texttt{plan}, \texttt{recover}\}$ classify the operational
structure of each trajectory step.
Table~\ref{tab:gpqa-action-motifs} reports the top 2-gram and 3-gram action
motifs by harness and outcome.
A motif's rate is the mean number of occurrences per trajectory.

\begin{table}[t]
\centering
\small
\caption{%
  Top action motifs by harness and outcome on GPQA-Diamond.
  Rate = mean occurrences per trajectory.
  Motifs ordered by rate within each (harness, outcome) cell.
}
\label{tab:gpqa-action-motifs}
\setlength{\tabcolsep}{3pt}
\begin{tabular}{llllr}
\toprule
Harness & Outcome & \multicolumn{2}{c}{Top 3 motifs} & Rate \\
\midrule
\multirow{3}{*}{OpenClaw}
  & \multirow{3}{*}{C}
    & answer$\to$verify            & & 0.53 \\
  & & answer$\to$answer            & & 0.22 \\
  & & answer$\to$verify$\to$recover &  & 0.15 \\
\cmidrule{2-5}
  & \multirow{3}{*}{W}
    & answer$\to$verify            & & 0.45 \\
  & & answer$\to$plan              & (re-planning) & 0.27 \\
  & & recover$\to$recover          & (oscillation) & 0.19 \\
\midrule
\multirow{3}{*}{OpenCode}
  & \multirow{3}{*}{C}
    & answer$\to$reason            & & 0.86 \\
  & & reason$\to$verify            & & 0.49 \\
  & & tool\_use$\to$reason         & (tool-grounded) & 0.42 \\
\cmidrule{2-5}
  & \multirow{3}{*}{W}
    & answer$\to$reason            & & 0.92 \\
  & & reason$\to$verify            & & 0.46 \\
  & & reason$\to$plan              & (stalled planning) & 0.23 \\
\midrule
\multirow{3}{*}{ZeroClaw}
  & \multirow{3}{*}{C}
    & reason$\to$reason            & & 11.43 \\
  & & reason$\to$reason$\to$reason & & 9.77 \\
  & & plan$\to$reason              & & 1.23 \\
\cmidrule{2-5}
  & \multirow{3}{*}{W}
    & reason$\to$reason            & & 10.88 \\
  & & reason$\to$reason$\to$reason & & 9.03 \\
  & & reason$\to$plan              & (diverted planning) & 1.44 \\
\bottomrule
\end{tabular}
\end{table}

\findingbox{The dominant action motif on OpenClaw wrong trajectories shifts
from productive verification (answer$\to$verify) to re-planning without
execution (answer$\to$plan) and recovery oscillation (recover$\to$recover).
OpenCode wrong trajectories lose the tool-grounded reasoning cycle
(tool\_use$\to$reason) that anchors its correct trajectories.}

OpenClaw correct trajectories follow a verify-then-recover cycle:
answer$\to$verify at 0.53, with answer$\to$verify$\to$recover at 0.15.
On wrong trajectories, the second-rank motif becomes answer$\to$plan at 0.27:
the model stops verifying and starts re-planning without executing.
The third-rank motif recover$\to$recover at 0.19 is the action-level analogue
of the ARM-level PD$\to$PD sink: recovery chains into itself without resolution.

OpenCode correct trajectories are anchored by tool\_use$\to$reason at 0.42.
On wrong trajectories, this motif drops below the top 3, replaced by
reason$\to$plan at 0.23: the model substitutes planning for execution while
maintaining the reason$\to$verify cycle at a comparable rate (0.46 vs.\ 0.49).

ZeroClaw trajectories are dominated by reason$\to$reason chains at
approximately 10 per trajectory regardless of outcome, the template signature.
The only outcome-conditioned signal is a slight shift from plan$\to$reason
(correct, 1.23) to reason$\to$plan (wrong, 1.44): wrong trajectories interrupt
reasoning with planning.

% ---------------------------------------------------------------------------
\paragraph{Action-conditioned entropy: tool use suppresses but does not
guarantee.}
% ---------------------------------------------------------------------------

Table~\ref{tab:gpqa-action-entropy} reports mean per-token log-probability
entropy conditioned on the presence of specific action patterns.

\begin{table}[t]
\centering
\small
\caption{%
  Mean per-token entropy conditioned on action pattern, by harness and outcome
  on GPQA-Diamond.
  ``Has tool'' = trajectory contains $\ge$1 tool\_use event.
  ``Has recover'' = trajectory contains $\ge$1 recover event.
  $\Delta$Ent = entropy(wrong) $-$ entropy(correct).
}
\label{tab:gpqa-action-entropy}
\setlength{\tabcolsep}{3pt}
\begin{tabular}{llrrrr}
\toprule
Harness & Action pattern & Ent. (C) & Ent. (W) & $\Delta$Ent & $n$ (C\,/\,W) \\
\midrule
\multirow{4}{*}{OpenClaw}
  & Has tool use      & 0.216 & 0.264 & +0.048 & 29\,/\, 5 \\
  & No tool use       & 0.247 & 0.125 & $-$0.122 & 102\,/\,62 \\
  & Has recover       & 0.201 & 0.169 & $-$0.032 & 30\,/\, 7 \\
  & Verify after ans. & 0.242 & 0.148 & $-$0.094 & 74\,/\,39 \\
\midrule
\multirow{4}{*}{OpenCode}
  & Has tool use      & 0.251 & \textbf{0.378} & \textbf{+0.127} & 80\,/\, 8 \\
  & No tool use       & 0.307 & 0.273 & $-$0.034 & 65\,/\,45 \\
  & Has verify        & 0.293 & 0.280 & $-$0.013 & 86\,/\,26 \\
  & Has plan          & 0.289 & 0.265 & $-$0.024 & 44\,/\,16 \\
\midrule
\multirow{3}{*}{ZeroClaw}
  & Has verify        & 0.245 & 0.350 & \textbf{+0.105} & 90\,/\,21 \\
  & Has plan          & 0.239 & 0.352 & \textbf{+0.113} & 122\,/\,27 \\
  & Sustained reason  & 0.245 & 0.359 & \textbf{+0.114} & 147\,/\,34 \\
\bottomrule
\end{tabular}
\end{table}

\findingbox{Tool use is an entropy suppressor, not a correctness guarantee.
Under OpenCode, tool-using correct trajectories have lower entropy than
non-tool-using ones (0.25 vs.\ 0.31), but tool-using wrong trajectories spike
to the highest entropy in the table (0.38).
The model persists in tool use under uncertainty, and that uncertainty remains
visible in the output.}

Three regularities emerge from Table~\ref{tab:gpqa-action-entropy}.
First, tool use lowers entropy on correct trajectories across both tool-using
harnesses: OpenClaw correct with tools 0.22 versus 0.25 without; OpenCode
correct with tools 0.25 versus 0.31 without.
Tool execution anchors reasoning and the output becomes more deterministic.

Second, tool use on wrong trajectories carries the opposite signal.
OpenCode wrong trajectories with tools have entropy 0.38, the highest cell in
the table: the model invokes tools under uncertainty, the output does not
resolve it, and subsequent reasoning retains high entropy.
This is the failure mode the ARM analysis identifies as uncertainty without
resolution.

Third, ZeroClaw's action-conditioned entropy shows a uniform positive
$\Delta$Ent: wrong trajectories have higher entropy than correct ones for every
action pattern (+0.11 for verify, +0.11 for plan, +0.11 for sustained reason).
The model is uncertain, the entropy reflects that uncertainty, but the template
prevents the uncertainty from spiraling into repetition: the entropy stays high
rather than collapsing, consistent with Finding~5's observation that ZeroClaw
inverts the entropy pattern.
